\documentclass[12pt,a4paper]{article}

% Русский язык
\usepackage[T2A]{fontenc}
\usepackage[utf8]{inputenc}
\usepackage[russian]{babel}

% Математика и графика
\usepackage{amsmath,amssymb}
\usepackage{graphicx}
\usepackage{float}

% Таблицы
\usepackage{tabularx}
\usepackage{booktabs}
\usepackage{array}
\usepackage{caption}
\captionsetup[table]{justification=raggedleft,singlelinecheck=false}
\setlength{\tabcolsep}{4pt}

% Ссылки и оглавление
\usepackage[hidelinks]{hyperref}
\usepackage{tocloft}

% Листинги кода
\usepackage{listings}
\usepackage{fancyvrb}
\lstset{
    basicstyle=\ttfamily\small,
    breaklines=true,
    frame=single,
    literate=
      {—}{{--}}1
      {←}{{$\leftarrow$}}1
      {→}{{$\rightarrow$}}1
}

% Геометрия и отступы
\usepackage{geometry}
\geometry{left=2.5cm,right=2.5cm,top=2cm,bottom=2cm}
\setlength{\parindent}{1.25cm}
\setlength{\parskip}{0pt}

% Формат заголовков
\usepackage{titlesec}
\titleformat{\section}{\large\bfseries}{\thesection}{1em}{}
\titleformat{\subsection}{\normalsize\bfseries}{\thesubsection}{1em}{}

\begin{document}

% Титульный лист
\begin{titlepage}
    \centering
    {\large \textbf{САНКТ-ПЕТЕРБУРГСКИЙ ГОСУДАРСТВЕННЫЙ УНИВЕРСИТЕТ}} \\[0.5cm]
    {\large Факультет прикладной математики--процессов управления} \\[0.5cm]
    {\large Программа бакалавриата} \\[0.2cm]
    {\large \textbf{«Большие данные и распределенная цифровая платформа»}} \\[4.5cm]

    {\Large \textbf{ОТЧЕТ}} \\[0.3cm]
    {\large По лабораторной работе № 2} \\[0.2cm]
    {\large По дисциплине «Алгоритмы и структуры данных»} \\[0.2cm]
    {\large На тему «Исследование задачи коммивояжера через алгоритм имитации отжига и муравьиный алгоритм с модификациями»} \\[4cm]

    \begin{flushright}
        Студент гр. 24Б15-пу \\
        Телицин А.Д. \\[0.5cm]
        Преподаватель \\
        Дик А.Г. \\[1cm]
    \end{flushright}

    \vfill
    \centering
    Санкт-Петербург \\
    2026 г.
\end{titlepage}

\tableofcontents
\newpage

\section{Цель работы}
\label{sec:goal}

Целью работы является исследование и сравнительный анализ качества решения задачи коммивояжера с помощью:
\begin{enumerate}
    \item базового алгоритма имитации отжига;
    \item модифицированного алгоритма имитации отжига (больцмановское охлаждение);
    \item базового муравьиного алгоритма;
    \item модифицированного муравьиного алгоритма («начальное расположение»).
\end{enumerate}

Основной акцент в сравнении делается на качестве получаемого маршрута (длине гамильтонова цикла), затем на времени выполнения и стабильности результатов при изменении параметров. Эксперименты проводятся на среднеразмерном графе Berlin52 и крупном графе World666.

\end{document}
